\documentclass[UTF8,a4paper,11pt]{ctexart}
\usepackage[a4paper,margin=2.35cm]{geometry}
\usepackage{amsmath,amssymb,booktabs,longtable,array,enumitem,microtype}
\usepackage{hyperref}
\hypersetup{colorlinks=true,linkcolor=black,urlcolor=blue}
% 字体按名字探测：原稿写死 C:/Windows/Fonts 路径，在非 Windows 机器上编译直接失败。
% 退役条件：构建环境固定为单一字库（容器内预装 Noto Serif/Sans SC）后，删去分支只留该字库。
\IfFontExistsTF{Noto Serif SC}{\setCJKmainfont{Noto Serif SC}}{\setCJKmainfont{Songti SC}}
\IfFontExistsTF{Noto Sans SC}{\setCJKsansfont{Noto Sans SC}}{\setCJKsansfont{Hiragino Sans GB}}
\IfFontExistsTF{TeX Gyre Termes}{\setmainfont{TeX Gyre Termes}}{\setmainfont{Times New Roman}}
\setlength{\parindent}{2em}
\setlength{\parskip}{0.35em}
\setlist{nosep,leftmargin=2.5em}
\renewcommand{\arraystretch}{1.25}
\title{\textbf{当 Peer MAX 溢出：产品市场联系、投资者注意力与收益反转}}
\author{Yangyi Liu\\西南交通大学}
\date{研究计划（中文版）\\2026年7月13日}

\begin{document}
\maketitle

\begin{abstract}
本研究考察：一家企业实现异常高的日收益后，是否会影响那些自身并未实现该收益、但在经济上与其相关的企业。核心处理变量是源企业产品市场近邻的连续、网络加权 raw MAX 日收益（\emph{PeerMAXraw}）。研究首先检验 PeerMAXraw 是否在控制 focal 企业自身 MAX 和普通 peer 收益后，预测 focal 股票下个月较低的收益；随后利用同一批源企业 MAX 实现构造日度事件研究，考察注意力是否转向经济上相关的企业、这些企业是否先出现价格上涨，以及该上涨是否随后反转。预先确定的产品相似度决定潜在传导方向，事前的投资者可识别性则区分投资者更可能感知的经济联系与主要由研究者观察到的联系。实证设计把 MAX 作为观察到的处理变量，把 lottery-like payoff 作为理论解释，而不是先验假设。研究还将使用市场或行业调整的 MAX、波动率标准化 MAX、信息离散度、高阶矩、新闻、分析师修正和未来基本面等检验相邻解释。因此，项目的结论将保持条件性：只有当月度定价事实、直接注意力反应、focal 企业的初始上涨和随后的反转均沿着投资者可识别的产品市场联系出现时，证据才支持“收益显著性能够沿经济关联资产传播，并产生暂时性价格压力”的解释。
\end{abstract}

\noindent\textbf{关键词：}最大日收益；产品市场近邻；投资者注意力；收益可预测性；价格压力；收益反转

\section{研究动机与问题}
股票在一个月内实现的最大日收益是一个简单、且具有明确经济含义的统计量。Bali 等（2011）发现，自身 MAX 较高的股票随后获得较低的平均收益，这一结果通常与投资者对 lottery-like payoff 的需求有关。既有结果主要是自身股票层面的：股票实现极端正收益，吸引需求，并相对于其后续收益变得偏贵。然而，投资者并不会孤立地评价证券。企业被组织在产品、技术、行业、主题和竞争集合中。因此，一家企业的显著收益可能改变投资者分配给其他看起来相关企业的注意力。

本计划询问 MAX 是否会沿这些经济联系传播。设源企业 $j$ 实现异常高的日收益，而 focal 企业 $i$ 销售相似产品。源企业的收益可直接从市场数据观察到，也可能出现在收益排行榜、新闻推送或讨论平台中。认识到两家企业联系的投资者可能会把注意力转向 $i$，增加交易，并暂时推高其价格，尽管 $i$ 并未实现源企业的收益。如果该反应体现的是投机性类别泛化，而不是共同现金流信息，那么 focal 价格上涨应当随后反转。核心机制是：

\begin{center}
\textbf{源企业 MAX} $\longrightarrow$ \textbf{对可识别产品近邻的注意力} $\longrightarrow$ \textbf{focal 上涨} $\longrightarrow$ \textbf{反转}。
\end{center}

项目围绕三个问题展开。第一，连续的月度 peer MAX 暴露是否预测 focal 股票下个月较低的收益？这是总体资产定价事实，也是与 MAX 文献建立联系的必要条件。第二，在显著的源企业 MAX 事件之后，产品相似度更高的企业是否获得更多注意力，尤其是在这种联系于事前更容易被投资者识别时？这是传导问题。第三，随后的价格路径更符合暂时性价格压力，还是信息扩散？单独的负未来收益不足以证明反转。focal 股票必须先出现异常上涨，再回吐这部分涨幅。若价格持续上涨且类别范围内出现持续的基本面修正，则更接近信息扩散。

处理与解释的区分是设计的核心。实证处理变量是由公司行为调整后的 raw return 构造的 $PeerMAX^{raw}$。Lottery-like payoff 是提出预测的经济解释。月度处理变量不使用研究者选定的尾部阈值、标准化异常收益或信息离散度指标。研究要检验的是：透明的 MAX 统计量是否具有增量预测内容，以及其日度表现是否产生解释所要求的注意力与价格路径。

项目位于四类既有发现的交叉处：自身 MAX 能预测收益（Bali 等，2011）；注意力可以在展示的证券之间联动或溢出（Drake 等，2017；Chen 等，2023）；渐进式与离散式收益路径中包含的信息含量不同（Da 等，2014）；经济相关企业之间存在跨股票收益可预测性（Hoberg and Phillips，2018；Avramov 等，2026）。本项目试图把这些文献连接成一个结果：一家企业的可见正收益，沿着预先确定且投资者可识别的产品市场网络，预测其他企业的注意力、暂时性上涨和反转。

\section{概念框架与可检验预测}
\subsection{从源企业收益到 focal 股票}
框架从有限注意力和基于类别的搜索出发。源企业的大幅正收益比普通收益更显著，可能促使投资者搜索相关名称、比较同一产品类别中的企业，或把感知到的收益机会外推到看起来相似的企业。产品市场相似度决定哪些企业是这种搜索的经济上合理的目标；可识别性决定该联系是否足够容易被投资者看见并引导注意力。

这一框架不同于纯粹的共同信息解释。共享信息同样可能使产品市场近邻产生相关反应，但它预测价格以及最终基本面持续同向运动。暂时性投机压力则预测有序路径：focal 注意力上升，短期交易和收益增加，随后价格上涨反转。实证分析不会先验地把其中任何一种解释视为成立，而是使用时序、价格符号、信息诊断和可识别性的横截面差异加以区分。

\subsection{预测}
\textbf{预测 1：总体定价。} 月度高 peer MAX 暴露越大，focal 股票在 $t+1$ 月的收益越低：
\[
\frac{\partial E[R_{i,t+1}]}{\partial PeerMAX^{raw}_{i,t}}<0.
\]
该关系应在控制 focal 企业自身 MAX 和 peer 普通累计收益后仍然存在。这一定义总体资产定价事实，但本身不能识别注意力或错误定价。

\textbf{预测 2：网络传导。} 源企业 MAX 事件之后，focal 企业的注意力反应应随事前产品相似度增加而增强；当源—focal 关系事前更容易识别时，这一梯度应更强：
\[
\frac{\partial AttentionResponse_{i,j,\tau}}{\partial Similarity_{ij,y-1}}>0,
\qquad
\frac{\partial^2 AttentionResponse_{i,j,\tau}}
{\partial Similarity_{ij,y-1}\partial Recognizability_{ij,y-1}}>0.
\]
仅使用相似度可能捕捉共同经济暴露，因此与可识别性的交互项对于解释“投资者注意力路径”很重要。

\textbf{预测 3：有序价格压力。} 强且可识别的 focal 近邻应当表现为：差异化事前趋势平坦；$\tau+1$ 至 $\tau+2$ 注意力上升；$\tau+1$ 至 $\tau+3$ 异常交易和 focal 累计异常收益为正；$\tau+4$ 至 $\tau+20$ 的增量收益为负。只有观察到初始上涨，才使用“反转”一词。没有早期上涨而只有较低未来收益，可能来自遗漏风险、延迟信息或一般可预测性，而不是事件诱发的价格压力反转。

\textbf{预测 4：信息含量。} 当源事件的观测信息含量更高时，暂时性压力模式应减弱。若事件之后出现分析师修正、持续的源企业和类别收益，或对应的经营与盈利实现，则应有更多延续、更少反转。相反，若注意力和 focal 上涨很强，但后续基本面修正很弱，则更符合暂时性压力。没有任何单一新闻过滤器能够识别“非信息冲击”。

\textbf{预测 5：收益方向与源企业聚合。} 正负极端收益都可能吸引注意力，但如果投资者是在外推上涨机会，正 MAX 更可能引起 focal 买入和正向上涨。若反应主要由最显著的单一 peer（Top1）驱动，更接近显著样本企业解释；若反应由大量 peer 的上涨广度驱动，则更接近类别范围的信息解释。这些是细化解释的证伪检验，不替代预测 1--3。

\section{数据与变量构造}
\subsection{样本与数据架构}
样本拟包括拥有日度和月度市场数据的中国 A 股普通股，以测量收益、涨跌停、交易活动和标准企业特征。会计数据用于规模、账面市值比和其他控制变量。近期 IPO、微型市值股票以及收益测量或交易限制使解释困难的观测将被排除或单独分析。结果同时报告等权和市值加权版本，避免总体事实完全由规模小、难交易的股票驱动。

产品网络采用由年报业务描述构造并预验证的中国 TNIC-inspired 文本相似度网络：同文本年的公司两两计算余弦相似度，用平均链聚类划出对手集合，团数逐年校准至目标中位 peer 数。该构造是传递划分，与 Hoberg--Phillips 的非传递阈值图（以相似度分布分位数为边）不同，故不写作 TNIC 复现；校准目标值取自申万二级中位 peer 的外部锚，其可复算值与外部锚的差异在参数登记中列明。为防止前视偏差和事件诱发的联系形成，日历年 $y$ 的主网络至少使用一个完整财政年度以前可获得的业务描述。注意力数据在股票—日层面测量，包括百度对股票名称或代码的搜索；在覆盖允许时，加入东方财富股吧的阅读、发帖或评论活动。信息变量结合交易所与巨潮资讯披露、新闻、分析师预测修正以及未来基本面结果。交易反应包括异常换手率和交易笔数；若交易数据允许，再补充订单流变量。

第一阶段是数据门，而不是回归。它将核验市场、网络、注意力和信息数据的共同覆盖；检查搜索关键词质量；统计可用源事件；检查事件窗口重叠；并确定是否能够形成强、弱和匹配比较企业。如果共同样本或关键词映射过于稀疏，无法可信地测量注意力机制，则不继续推进注意力分析。

\subsection{核心 MAX 处理变量}
源企业 $j$ 在月份 $t$ 的 raw MAX 定义为：
\begin{equation}
MAX^{raw}_{j,t}=\max_{d\in t}r^{raw}_{j,d},
\end{equation}
其中 $r^{raw}_{j,d}$ 是经过公司行为调整的日收益，不扣除市场或行业收益，也不除以估计波动率。令 $Similarity_{ij,y-1}$ 表示 focal $i$ 与源企业 $j$ 的事前产品相似度，在 focal 企业的 peer 集合内归一化：
\begin{equation}
w_{ij,y-1}=\frac{Similarity_{ij,y-1}}
{\sum_{k\ne i}Similarity_{ik,y-1}}.
\end{equation}
核心处理变量为：
\begin{equation}
PeerMAX^{raw}_{i,t}=\sum_{j\ne i}w_{ij,y-1}MAX^{raw}_{j,t}.
\end{equation}
这一连续变量直接连接自身 MAX 文献，不需要临界值 $c$，并贯穿月度分析。市场或行业调整 MAX、标准化异常收益 MAX 仅作为诊断：前者检验关系是否由市场或行业冲击驱动，后者检验 raw MAX 是否只是在选择波动率不同的企业。它们不替代式（3）。若只有研究者标准化的变量有效，而 raw MAX 无效，则很难维持“直接可观察收益”的解释。

\subsection{控制变量与竞争变量}
核心控制变量是 focal 企业自身的 $OwnMAX^{raw}_{i,t}$ 和网络加权的 peer 累计月度收益 $PeerReturn_{i,t}$。其他 peer 变量包括特质波动率（PeerIVOL）、偏度（PeerSkew）和信息离散度（PeerID）。这些变量都由日收益共同生成，不能构成结构性分解；其作用较窄，只是检验 PeerMAX 是否超越普通升值、波动率、高阶矩和收益路径形状。

\subsection{可识别性、注意力与信息}
投资者可识别性使用源事件之前的信息测量。优先使用与年报文本不同来源的事前新闻共同提及、论坛共同提及，以及可获得时的搜索或分析师报告关联。共享产品关键词可以补充，但会与文本网络分开，因为两者可能机械重叠。目标不是断言每个投资者都看见每条联系，而是检验经济联系越具有事前公共可见性，传导是否越强。

直接注意力结果包括异常百度搜索和异常股吧阅读、发帖或评论。交易结果包括异常换手率、交易笔数、订单到达强度，以及可获得时的主动买卖不平衡。小单不平衡最多是投机交易代理变量。搜索、平台使用、换手率和小额交易不能在没有独立验证的投资者类型数据时识别零售投资者，因此本计划不提出投资者身份结论。

信息含量通过多个互补指标诊断：事件附近的源企业公告与新闻、后续分析师预测修正、后续盈利和销售等经营结果、源企业收益持续性、类别范围的收益和基本面延续，以及 focal 企业自己的新闻或互动平台披露。文章将使用“观测到的信息较少”，而不使用“干净事件”或“非信息事件”等表述。

\section{月度资产定价设计}
\subsection{基准检验}
核心月度规格为 Fama--MacBeth 横截面回归：
\begin{equation}
R_{i,t+1}=\alpha_t+\beta PeerMAX^{raw}_{i,t}+\theta OwnMAX^{raw}_{i,t}
+\phi PeerReturn_{i,t}+\Gamma'X_{i,t}+\varepsilon_{i,t+1},
\end{equation}
其中 $X_{i,t}$ 包括规模、价值、动量、短期反转、流动性、换手率、分析师覆盖和网络度数。主要预测为 $\beta<0$。报告系数的时间序列平均值、PeerMAX 四分位距或标准差变化对应的经济意义，以及适合收益期限的 Newey--West 推断。

组合检验按 PeerMAXraw 对股票排序，同时报告等权和市值加权收益。两个加权口径的并报用于诊断结论对规模权重的敏感性：若两者符号分裂且均不显著，按主效应未检出处理，不据此主张任一方向上的定价事实；尤其是市值加权口径功效较低时，其不显著不能读作该口径下无效应。双重排序或特征调整后的组合收益用来判断收益差是否在自身 MAX 和普通 peer 收益之后仍然存在。还将在排除微型股、近期 IPO 和受涨跌停影响的源观测后评估结果。这些限制不是为了寻找最佳子样本，而是检验定价事实是否超越最难投资的股票仍具经济意义。

\subsection{增量内容与函数形式}
独特性分析采用嵌套且透明的规格，按顺序加入自身 MAX、PeerReturn、PeerIVOL 与 PeerSkew，以及 PeerID。这一顺序询问 PeerMAX 是否超越 focal 自身 lottery demand、普通 peer 升值、波动或偏度环境以及收益路径离散度。控制变量加入后仍然存在，只能说明增量预测内容，并不等于纯结构性 MAX 效应，也不等于投资者偏好本身。若 PeerMAX 被相邻变量完全吸收，独立 MAX 传导主张即失败，应改写为更一般的 peer 可预测性项目。在 PeerIVOL、PeerSkew 与 PeerID 三列实现之前，独特性检验只覆盖 PeerReturn；此阶段独立 MAX 传导主张记为未检验，不得写成已通过。

月度核心变量保持连续形式。为避免施加任意临界值，将使用 PeerMAX 分组和限制性三次样条检查函数形式。正尾部成分 $(MAX-c)^+$、Top1 peer 暴露和 peer jackpot 的广度只作为连续关系的分解。平滑关系意味着不需要尾部阈值；上尾部集中支持显著 jackpot 概括；Top1 占优强化样本企业注意力解释；广度占优强化类别信息解释。

月度与日度部分将在源企业实现层面连接。事件研究不能因为某些极端观测产生漂亮事件图，就另行挑选一组不相关的观测。对月度 PeerMAX 有贡献的源企业 MAX 观测将保留标识，使日度分析能够聚合或分解同一暴露。只有保持这一一致性，事件证据才可能解释月度定价事实。

\section{源事件与事件时间设计}
\subsection{事件定义}
月度定价检验不需要阈值，但日度事件分析需要规则来确定显著源企业实现。对源企业 $j$ 在日期 $\tau$ 的事件定义为：
\begin{equation}
Jackpot_{j,\tau}(q)=1\left\{Rank^{board\times day}_{j,\tau}\ge q\right\},
\end{equation}
其中排名依据相关交易板块和日期内、经公司行为调整的 raw 日收益。主展示使用 $q=0.99$，取一个事件序列中的首次越过，并对 $q\in\{0.95,0.975,0.99,0.995\}$ 报告反应路径。该网格必须在查看 focal 注意力或未来收益前选定。

该规则不使论文的主要结果依赖阈值；它的唯一作用是为事件研究确定日期。按板块—日期排名可以适应不同涨跌停制度。由于这种排名即使在平静日也可能选出最高百分位，研究将报告源企业绝对收益分布并核验事件是否具有经济显著性。涨跌停事件将从主要非涨跌停样本排除并单独分析：它们可能是最可见的 A 股事件，但收益受到交易规则删失，价格路径需要单独处理。

\subsection{比较组与回归设计}
对每个源事件 $(j,\tau)$，focal 企业分为三个具有经济含义的组：产品相似度高的强 peer；同一大行业但产品相似度低的弱 peer；以及事前特征相似、但没有产品联系的匹配非 peer。事件图展示三组路径以及强 peer 减弱 peer、或强 peer 减匹配非 peer 的差异。连续规格使用完整相似度梯度。

对事件时间跨度 $h$ 内的结果 $Y$，基准堆叠规格为：
\begin{equation}
Y_{i,j,\tau,h}=\alpha_{j,\tau,h}+\mu_i+\beta_h Similarity_{ij,y-1}
+\rho_h Similarity_{ij,y-1}\times Recognizability_{ij,y-1}
+\Gamma'_hX_{i,\tau-1}+\varepsilon_{i,j,\tau,h}.
\end{equation}
源事件—时间跨度固定效应保证 focal 企业是在同一源企业实现条件下比较；focal 企业固定效应吸收其持续性特征。交互项系数 $\rho_h$ 检验当经济联系在事件前更可识别时，传导是否更强。控制变量和匹配变量均需早于 $\tau$。

推断必须反映源企业重复出现、focal 企业重复出现和共同日历日期。主实现将在实际样本结构允许时采用多维聚类，并报告源事件层面和 focal 企业层面的聚类敏感性。保持度数和行业结构的网络重连将生成安慰剂分布，在保留重要网络结构的同时打断真实 source--focal 联系。该随机化检验补充而不替代常规推断。

重叠事件需要特别处理。不能把一个 focal 日期的结果复制到多个同时发生的源事件中，再声称通过 focal-date 固定效应识别。分析将按预先确定的规则删除重叠窗口，或在 focal-date 层面聚合同一日期的同时源暴露。具有自身公告、重大新闻或互动平台披露的 focal 观测，将从最干净的事件比较中删除并单独分析。

\subsection{有序结果}
事件分析是一条有序序列，而不是一组彼此无关的显著系数。差异化事前趋势必须平坦；focal 搜索或平台注意力应在 $\tau+1$ 至 $\tau+2$ 上升；异常交易和 focal 累计异常收益应在 $\tau+1$ 至 $\tau+3$ 为正；之后 $\tau+4$ 至 $\tau+20$ 的增量 focal 收益应为负，并反转早期上涨。

focal 异常收益将相对于市场、行业和匹配企业基准报告，但最有信息量的图是同一源事件内强 peer 与弱 peer 或匹配 peer 的差异。该差异更透明地吸收市场条件和共同源企业实现。没有日内阈值穿越时间时，不能确定源收益、focal 注意力和 focal 收益在同一天内的先后。因此基准因果顺序从 $\tau+1$ 开始。若以后获得日内数据，严格的日内排序属于额外模块，而不是基准设计的必要条件。

注意力、交易和收益是共同内生的响应。不能通过在收益回归中加入注意力变量后观察系数衰减，来推断因果中介。研究将把有序时序、预先确定的网络梯度、可识别性交互和后续价格反转作为相互强化、但与机制一致的证据。

\section{替代解释与证伪}
\subsection{信息扩散}
首要替代解释是：源企业 MAX 反映了与产品需求、技术、监管或投入市场有关的信息，并且这些信息对相似企业确实相关。在这种情况下，产品相似度正确地预测 focal 反应，而不是因为注意力导致错误定价。源事件固定效应并不能消除这一解释，因为相似度更高的 focal 企业可能对同一新闻具有更强的基本面暴露。

因此研究将结合公开信息分类与后续结果。公告和新闻指标描述观测到的公开信息，但不能识别外生或非信息冲击。分析师修正、未来盈利、销售和经营结果检验事件后是否出现新的现金流信息。源收益持续性和类别范围的延续检验反应是否持久。若可识别 peer 出现注意力和初始上涨，而后续基本面修正很弱且价格上涨反转，则支持暂时性压力；若价格和基本面反应同向持续，则支持信息扩散。

\subsection{收益路径、风险与高阶矩}
PeerReturn、PeerIVOL、PeerSkew 和 PeerID 进入月度独特性检验，因为大 MAX 与累计收益、波动率、偏度以及正负收益日的分布存在机械关系。控制这些变量不能声称得到纯结构性 MAX 处理；研究只能提出一个更可防守的问题：MAX 统计量是否含有增量信息。

\subsection{次级证伪检验}
在最低证据包成功后，再比较正 MAX 与绝对收益匹配的负极端收益；两者都可能增加注意力，但正事件之后更强的 focal 买入和上涨有助于区分上涨机会吸引与一般极端事件显著性。Top1 与广度分解区分显著样本企业和类别信息。涨跌停与非涨跌停事件分开报告。行业、共同分析师、地理位置、公司名称、证券代码邻近度和随机网络等替代网络，用来检验产品联系是否超越机械或制度性邻近。弱联系企业、随机日期和事前伪事件提供安慰剂路径。

这些检验都不能替代核心序列。如果 PeerMAX 没有独立的月度定价事实、focal 注意力没有响应、focal 价格没有先上涨，或初始上涨没有反转，再长的稳健性附录也不能挽救项目。

\section{贡献、可行性与决策规则}
若完整证据链成立，研究有三项相关贡献。第一，将自身 MAX 文献扩展到：没有实现该收益、但经济上相关的资产也会受到极端正收益影响。第二，刻画一个注意力网络：当产品关系对投资者更可识别时，预先确定的产品相似度传导最强。第三，把注意力传导与有序价格路径——先上涨、后反转——联系起来，而不是把较低未来收益直接视为错误定价。

主张将受设计边界约束。源企业 MAX 事件并非外生；产品相似度梯度本身不能证明注意力传导；没有观察到新闻不等于没有信息；搜索、换手、交易规模和订单不平衡不能识别零售投资者。因此研究对象是投资者注意力和暂时性交易压力，而不是某一种投资者类型。设计提供的是严谨的事件内比较和一致的结果序列，而不是因果中介估计。

项目按决策门推进。G0 核验数据覆盖、关键词质量、源事件数量、比较组可行性和可识别性。M1 产出月度 PeerMAX 组合与回归证据，并配套样条或分组图。M2 检验关系能否在 peer 收益、风险、高阶矩和信息离散度之后仍然存在。M3 形成一个规模较小但完全连接的事件样本，产出注意力和 focal 累计异常收益图。只有当 M1--M3 足够有希望时，才进行完整信息、替代网络和交易数据合并。

停止规则与经济主张对应。若数据门失败，注意力项目停止；若 PeerMAX 不能预测 focal 收益，不存在独立的 MAX 定价论文；若 PeerMAX 被 peer 收益、风险或 ID 吸收，独特 MAX 渠道失败，项目应改写或并入更一般的 peer 可预测性研究；若定价事实成立但注意力不响应，必须删除行为注意力机制；若注意力和上涨发生但没有反转，则证据更接近共同注意力或信息扩散，必须放弃暂时性价格压力主张。只有完整的月度到日度链条，才足以支持以 MAX 相关暂时性价格压力为题的论文。

若整个序列在经济意义上足够清晰，并且能够通过可识别网络和信息替代解释，项目的自然目标是较强的《Journal of Empirical Finance》论文。若要达到《Journal of Banking \& Finance》的上行目标，则需要特别强的时间证据、专有注意力或订单流数据，或直接证明投资者确实感知产品联系。仅有一个一般性的 PeerMAX 系数，再加上换手率、小单或注意力异质性，达不到这一标准。

\section{最低证据包}
最低论文包由少量可决策的产出组成，而不是大量规格菜单。

\begin{itemize}
\item \textbf{表 1：}月度 PeerMAX 组合排序与 Fama--MacBeth 回归；必须得到稳定的负向下月收益关系。
\item \textbf{表 2：}嵌套独特性检验；必须表明关系没有被普通 peer 收益、风险、高阶矩或信息离散度吸收。
\item \textbf{图 1：}按相似度和可识别性展示事件时间注意力；必须有平坦的差异化事前趋势和事件后的注意力梯度。
\item \textbf{图 2：}同一源事件内 focal 异常收益差；必须展示先上涨、后反转。
\item \textbf{表 3：}网络与信息诊断；可识别联系上的传导应更强，而信息含量较高时反转应更弱。
\end{itemize}

只有这组证据成立后，才加入替代网络、更多异质性、阈值变体和精细交易指标。该顺序避免次级规格掩盖核心论证中缺失的环节。最终标准是：MAX 是处理变量，lottery-like payoff 是解释，而注意力反应和有序价格路径共同决定该解释是否可信。

\appendix
\section{变量层级与详细定义}
\begin{longtable}{p{0.18\textwidth}p{0.27\textwidth}p{0.43\textwidth}}
\toprule
角色 & 变量 & 目的\\
\midrule
核心处理 & $PeerMAX^{raw}$ & 网络加权、直接可观察的源企业收益；唯一的核心预测变量。\\
核心控制 & $OwnMAX^{raw}$、$PeerReturn$ & focal 自身 MAX 需求与普通 peer 升值。\\
独特性 & $PeerIVOL$、$PeerSkew$、$PeerID$ & 波动率、高阶矩和收益路径替代解释。\\
调整诊断 & $PeerMAXAR$、$PeerMAXZAR$ & 共同市场或行业冲击，以及波动率缩放。\\
函数分解 & $PeerTail$、Top1、$PeerShare$ & 上尾、显著样本企业和广度分解。\\
事件时间 & $Jackpot(q)$ & 显著源事件日期；不定义月度结果。\\
\bottomrule
\end{longtable}

所有 peer 变量均使用共同的预先确定产品网络权重，除非检验明确比较替代网络。没有任何变换或阈值化变量替代 raw PeerMAX 的核心规格。

\subsection{信息离散度}
沿用 Da 等（2014）的零收益修正源企业指标：
\begin{equation}
ID_{j,t}=sign(PRET_{j,t})\frac{\%Neg_{j,t}-\%Pos_{j,t}}
{\%Neg_{j,t}+\%Pos_{j,t}},\qquad
PeerID_{i,t}=\sum_jw_{ij,y-1}ID_{j,t}.
\end{equation}
报告一个月、三个月和十二个月窗口。一个月版本预计噪声较大，不作为决定性控制变量。ID 是竞争性的收益路径代理，不是直接注意力指标。

\section{详细经验规格}
\subsection{月度规格}
完整独特性规格为：
\begin{equation}
\begin{aligned}
R_{i,t+1}={}&\alpha_t+\beta PeerMAX^{raw}_{i,t}+\theta OwnMAX^{raw}_{i,t}
+\phi PeerReturn_{i,t}\\
&+\lambda_1PeerIVOL_{i,t}+\lambda_2PeerSkew_{i,t}
+\lambda_3PeerID_{i,t}+\Gamma'X_{i,t}+\varepsilon_{i,t+1}.
\end{aligned}
\end{equation}
函数形式规格以分组和限制性三次样条估计 $f(PeerMAX^{raw}_{i,t})$。组合排序、连续回归和调整收益定义将联合解释，而不是选择统计量最大的规格。

\subsection{回程诊断}
在累计收益单元格内，可选的源企业层面诊断规格为：
\begin{equation}
Y_{j,t}=\alpha+\beta MAX^{raw}_{j,t}+\phi MonthlyReturn_{j,t}
+\delta ID_{j,t}+\eta IVOL_{j,t}+\kappa Skew_{j,t}+\Gamma'Z_{j,t}+u_{j,t}.
\end{equation}
事件图可在累计收益、大行业、规模、换手率和观测新闻状态匹配的前提下，比较高 MAX 与低 MAX 的源月份。收益生成的变量之间不强制完全平衡；分析将报告共同支撑、剔除观测与回归调整后的估计。该诊断不是主检验、不是 lottery 偏好的识别策略，也不是停止规则。

\subsection{事件推断}
式（6）分别对注意力、交易和收益期限估计。推断将处理源企业和 focal 企业重复出现以及共同日期依赖。具体多维聚类根据实际依赖结构选择并透明报告。网络随机化保持源企业和 focal 企业度数以及大行业构成。重叠事件规则将在查看事件后结果之前冻结。

\section{测量与实施协议}
对每一个注意力序列，异常活动相对于排除即时事件窗口的事前滚动基准测量。关键词映射需人工审查含义模糊的公司名称和代码变化。平台结果先分别展示，再决定是否构造标准化指数。交易反应优先使用覆盖广的换手率和交易笔数；只有交易数据可靠时才加入订单不平衡和订单到达强度。

信息变量保持互补，不压缩成不透明的综合分数。设计区分源企业公开信息、focal 自身信息、后续分析师修正、后续实际基本面和价格持续性。某事件可能公开新闻较少但后续信息实现较高；这种情况不标记为“非信息事件”。

在估计前报告：各年份和板块覆盖；市场、网络和注意力样本交集；源事件数量；源企业绝对收益分布；每个源事件的强、弱和匹配 focal 企业数量；事件重叠率；可识别性覆盖；以及 focal 自身新闻排除频率。这些检查在大量数据构造继续前决定比较设计是否可行。

\section{快速失败工作计划}
\begin{center}
\begin{tabular}{p{0.12\textwidth}p{0.58\textwidth}p{0.18\textwidth}}
\toprule
阶段 & 可决策产出 & 预计预算\\
\midrule
G0 & 覆盖、关键词审计、事件数量、比较组可行性和可识别性覆盖 & 2--3 个 RA 日\\
M1 & PeerMAX 定价表、组合图与样条/分组图 & 3--5 个 RA 日\\
M2 & 包含 peer 收益、风险、高阶矩和 ID 的独特性表 & 2--4 个 RA 日\\
M3 & 带注意力和 focal CAR 图的连接事件样本 & 4--6 个 RA 日\\
扩展 & 完整信息、替代网络和交易数据合并 & 仅在 M1--M3 后\\
\bottomrule
\end{tabular}
\end{center}

每个阶段都必须以可复现代码、冻结的样本定义和一张可决策表或图结束。不能仅因为早期系数统计显著，就投入昂贵的数据合并；进入下一阶段必须同时满足经济意义上的符号和上文规定的路径。

\section*{参考文献}
\begin{itemize}
\item Ali, U., and D. Hirshleifer. 2020. Shared analyst coverage: Unifying momentum spillover effects. \emph{Journal of Financial Economics} 136, 649--675.
\item Avramov, D., S. Ge, S. Li, and O. Linton. 2026. Dual peer effects and cross-stock predictability. \emph{Journal of Financial Economics} 180, 104274.
\item Bali, T. G., N. Cakici, and R. F. Whitelaw. 2011. Maxing out: Stocks as lotteries and the cross-section of expected returns. \emph{Journal of Financial Economics} 99, 427--446.
\item Chen, X., L. An, Z. Wang, and J. Yu. 2023. Attention spillover in asset pricing. \emph{Journal of Finance} 78, 3515--3559.
\item Da, Z., U. G. Gurun, and M. Warachka. 2014. Frog in the pan: Continuous information and momentum. \emph{Review of Financial Studies} 27, 2171--2218.
\item Drake, M. S., J. Jennings, D. T. Roulstone, and J. R. Thornock. 2017. The comovement of investor attention. \emph{Management Science} 63, 2847--2867.
\item Hoberg, G., and G. Phillips. 2018. Text-based industry momentum. \emph{Journal of Financial and Quantitative Analysis} 53, 2355--2388.
\item Nguyen, H., and C. Truong. 2018. When are extreme daily returns not lottery? At earnings announcements! \emph{Review of Accounting Studies} 23, 1534--1562.
\end{itemize}

\end{document}
